
% !TEX root = allinone_main_v2.tex

\begin{thebibliography}{R-WW17}

\bibitem[Arl90]{Arlettaz} D. Arlettaz. The first~$k$-invariant of a double loop space is trivial. Arch. Math. 54 (1990) 84--92.

\bibitem[BP72]{BP72}M. Barratt, S. Priddy. On the homology of non-connected monoids and their associated groups. Comment. Math. Helv. 47 (1972) 1--14.

\bibitem[Bro87]{Bro87} K.S. Brown. Finiteness properties of groups. J. Pure Appl. Algebra 44 (1987) 45--75.

\bibitem[Bro92]{Brown:Suggestion} K.S. Brown. The geometry of finitely presented infinite simple groups. Algorithms and classification in combinatorial group theory (Berkeley, CA, 1989) 121--136. Math. Sci. Res. Inst. Publ. 23. Springer, New York, 1992.

\bibitem[BG84]{Brown+Geoghegan} K.S. Brown, R. Geoghegan. An infinite-dimensional torsion-free~$\rmF\rmP_\infty$ group. Invent. Math. 77 (1984) 367--381.

\bibitem[ER-W17]{EbertRW} J. Ebert, O. Randal-Williams. Semi-simplicial spaces. Preprint (2017). ArXiv:1705.03774.   

\bibitem[Far05]{Farley} D.S. Farley. Homological and finiteness properties of picture groups. Trans. Amer. Math. Soc. 357 (2005) 3567--3584.

\bibitem[GS87]{Ghys+Sergiescu} \'E. Ghys, V. Sergiescu. Sur un groupe remarquable de diff\'eo\-mor\-phismes du cercle. Comment. Math. Helv. 62 (1987) 185--239.

\bibitem[Gra76]{Grayson:QuillenII} D. Grayson. Higher algebraic K-theory. II (after Daniel Quillen). Algebraic K-theory (Proc. Conf., Northwestern Univ., Evanston, Ill., 1976) 217--240. Lecture Notes in Math. 551. Springer, Berlin, 1976.

\bibitem[HW10]{Hatcher+Wahl} A. Hatcher, N. Wahl. Stabilization for mapping class groups of 3-mani\-folds. Duke Math. J. 155 (2010) 205--269.

\bibitem[Hig74]{Higman} G. Higman. Finitely presented infinite simple groups. Notes on Pure Mathematics 8 (1974). Department of Pure Mathematics, Department of Mathematics, I.A.S. Australian National University, Canberra, 1974.

\bibitem[Kap02]{Kapoudjian} C. Kapoudjian. Virasoro-type extensions for the Higman--Thompson and Neretin groups. Q. J. Math. 53 (2002) 295--317.

\bibitem[McDS76]{McDuff-Segal} D. McDuff, G. Segal. Homology fibrations and the ``group-completion'' theorem. Invent. Math. 31 (1976) 279--284.  

\bibitem[Muk66]{Mukai} J. Mukai. Stable homotopy of some elementary complexes. Mem. Fac. Sci. Kyushu Univ. Ser. A 20 (1966) 266--282.

\bibitem[Qui73]{Quillen} D. Quillen. Higher algebraic K-theory. I. Algebraic K-theory, I: Higher K-theories (Proc. Conf., Battelle Memorial Inst., Seattle, Wash., 1972) 85--147. Lecture Notes in Math. 341. Springer, Berlin 1973. 

\bibitem[R-W13]{Randal-Williams} O. Randal-Williams. `Group-completion', local coefficient systems and perfection. Q. J. Math. 64 (2013) 795--803.

\bibitem[R-WW17]{Randal-Williams+Wahl} O. Randal-Williams, N. Wahl. Homological stability for automorphism groups. Adv. Math. 318 (2017) 534--626.

%\bibitem[Tho79]{Thomason1979} R.W.~Thomason. First quadrant spectral sequences in algebraic K-theory. Algebraic topology, Aarhus 1978 (Proc. Sympos., Univ. Aarhus, Aarhus, 1978) 332--355. Lecture Notes in Math. 763. Springer, Berlin, 1979.

\bibitem[Tho82]{Thomason} R.W.~Thomason. First quadrant spectral sequences in algebraic K-theory via homotopy colimits. Comm. Algebra 10 (1982) 1589--1668.

\bibitem[Thu17]{Thumann} W.~Thumann. Operad groups and their finiteness properties. Adv. Math. 307 (2017) 417--487. 

\bibitem[Zee64]{Zeeman} E.C.~Zeeman. Relative simplicial approximation. Proc. Cambridge Philos. Soc. 60 (1964) 39--43. 

\end{thebibliography}
